%       --- Constraints and Standards --- Constraints and Standards --- Constraints and Standards --- Constraints and Standards --- Constraints and Standards
\chapter{Constraints and Standards} % --- Constraints and Standards --- Constraints and Standards --- Constraints and Standards --- Constraints and Standards
%       --- Constraints and Standards --- Constraints and Standards --- Constraints and Standards --- Constraints and Standards --- Constraints and Standards

%       --- Constraints --- Constraints --- Constraints --- Constraints --- Constraints --- Constraints --- Constraints --- Constraints --- Constraints ---
\section{Constraints} % --- Constraints --- Constraints --- Constraints --- Constraints --- Constraints --- Constraints --- Constraints --- Constraints ---
%       --- Constraints --- Constraints --- Constraints --- Constraints --- Constraints --- Constraints --- Constraints --- Constraints --- Constraints ---

\begin{tcolorbox}[      % Delete me later
    boxrule=0pt,        % Delete me later
    colback=white,      % Delete me later
    coltext=red         % Delete me later
]                       % Delete me later
    Constraints are factors that limit your development process. They could be material such as money or hardware, or non-material such as time or customer satisfaction. Constraints include budget, customer satisfaction, quality, requirements, risk, time, ease of use, reliability, maintenance costs, etc.
    \vspace{15pt}
    
    Generally, a trade-off exists between time, cost, and scope (features and functionality) and should be discussed. Resources constraints include equipment, materials, software, facilities, and people. In addition, there are legal and ethical constraints to be considered when developing any project.
    \vspace{15pt}
    
    For example, will your project be maintained after deployment? Did the use of third-party software or hardware affect the design or your project? Did legal requirements such as HIPAA limit the scope of your project? Did the licensing of any components used affect your final project?
    \vspace{15pt}

    \textbf{HAVE AT LEAST 5 CONSTRAINTS!}
    \vspace{15pt}
    
    \textbf{REFERENCE AT LEAST ONE CONSTRAINTS YOU MENTION HERE!}

    \vspace{15pt}       % Delete me later
    \begin{center}      % Delete me later
    \textbf{Don't forget to delete this explanatory red text/block introduction, and replace it with your text in black.}
    \end{center}        % Delete me later
\end{tcolorbox}         % Delete me later

\newpage

%       --- Standards --- Standards --- Standards --- Standards --- Standards --- Standards --- Standards --- Standards --- Standards --- Standards ---
\section{Standards} % --- Standards --- Standards --- Standards --- Standards --- Standards --- Standards --- Standards --- Standards --- Standards ---
%       --- Standards --- Standards --- Standards --- Standards --- Standards --- Standards --- Standards --- Standards --- Standards --- Standards ---

\begin{tcolorbox}[      % Delete me later
    boxrule=0pt,        % Delete me later
    colback=white,      % Delete me later
    coltext=red         % Delete me later
]                       % Delete me later
    Standards are specifications or requirements, practices, guides, test methods, etc., covering various materials, products, systems or services. Standards are minimally accepted professional practices and qualities that must be observed. Common examples of standards in computer science and engineering include IEEE standards for electrical engineering and networking, and ISO standards for programming languages and design methodologies.
    \vspace{15pt}
    
    For example, IEEE provides standards for Ethernet (IEEE 802.3), wireless communications (IEEE 802.11), Bluetooth (IEEE 802.15), Verilog (IEEE 1364), as well as public-key cryptography (IEEE P1363). ISO provides standards for C++ (ISO 14882), UML (ISO 19501), QR codes (ISO 18004), ECMAScript (i.e., JavaScript) (ISO 16262), and software testing (ISO 29119-1), among others.
    \vspace{15pt}
    
    For example, did you choose specific standards for interoperability or to minimize power consumption? Did you use standard languages and libraries to help ensure maximum portability? Did your user interface conform to industry design standards for the devices used? Did you use any accessibility libraries or tools to make your project more inclusive?
    \vspace{15pt}

    \textbf{HAVE AT LEAST 5 STANDARDS!}
    \vspace{15pt}
    
    \textbf{YOU NEED TO REFERENCE EVERY STANDARD YOU MENTION HERE!}

    \vspace{15pt}       % Delete me later
    \begin{center}      % Delete me later
    \textbf{Don't forget to delete this explanatory red text/block introduction, and replace it with your text in black.}
    \end{center}        % Delete me later
\end{tcolorbox}         % Delete me later